% =============================================================================
% Major Technical Project report
% Extreme Value-Guided Deep Generative Models for Imbalanced Classification
% Case study: next-day temperature extremes in Himachal Pradesh
%
% Compile: pdflatex report.tex (twice). Figures are optional: put the PNGs from
% Drive (results/step3b/evt_qq_std.png, results/step4/tail_survival.png,
% results/step4/pr_curves.png) in a folder "figures/" next to this file.
% Values marked \pending{} are filled in after the Step 4 run.
% =============================================================================
\documentclass[11pt,a4paper]{article}
\usepackage[margin=2.3cm]{geometry}
\usepackage[T1]{fontenc}
\usepackage{lmodern}
\usepackage{amsmath,amssymb}
\usepackage{booktabs}
\usepackage{graphicx}
\usepackage{xcolor}
\usepackage{hyperref}
\usepackage{caption}
\usepackage{enumitem}
\usepackage{siunitx}
\sisetup{detect-all}
\setlist{nosep,leftmargin=1.4em}
\hypersetup{colorlinks=true,linkcolor=blue!50!black,citecolor=blue!50!black,urlcolor=blue!50!black}
\newcommand{\pending}{\textcolor{red!70!black}{\textit{pending}}}
\newcommand{\optfig}[2]{%
  \IfFileExists{#1}{\includegraphics[width=#2]{#1}}%
  {\fbox{\parbox{0.8\linewidth}{\centering\small Figure file \texttt{\detokenize{#1}} not found.\\
   Copy it from Google Drive into \texttt{figures/}.}}}}

\title{\textbf{Extreme Value-Guided Deep Generative Models\\for Imbalanced Classification}\\[4pt]
\large Case study: next-day temperature extremes in Himachal Pradesh, 1981--2025}
\author{Divyansh Jain\\Major Technical Project}
\date{\today}

\begin{document}
\maketitle

\begin{abstract}
Minority classes in imbalanced classification are often the most extreme cases in the data, yet the
standard remedies (class weighting, SMOTE) and modern deep generative models do not model tail
behaviour explicitly. We treat the minority class as a tail phenomenon: a Generalized Pareto
Distribution (GPD) fitted by peaks-over-threshold describes how severe minority examples are, and this
description constrains the synthetic examples produced by a conditional tabular GAN (CTGAN). We study
next-day extreme temperature days on a \ang{0.1} ERA5-Land grid over Himachal Pradesh (529 cells,
1981--2025, about 1\% positives). We build season-wise, whole-year and day-of-year extreme labels from
1981--2010 thresholds, show that standardising excesses by each cell's own variability turns
spuriously heavy pooled tails into bounded ones ($\xi\approx-0.15$ for whole-year heat), and establish
baselines on the unseen decade 2016--2025: adding threshold-distance features raises LightGBM's PR-AUC
from 0.43--0.45 to 0.53--0.56, with SMOTE best (0.562) but a persistence rule still ahead overall
(0.586). The data also show a doubling of heat extremes in the test decade and more intense winter
warm extremes. The comparison of vanilla and EVT-guided CTGAN is \pending.
\end{abstract}

\tableofcontents
\bigskip

% -----------------------------------------------------------------------------
\section{Introduction}
In imbalanced classification the class of interest is rare and is therefore learned poorly. Common
remedies re-weight the loss or create synthetic minority samples by interpolation
(SMOTE~\cite{chawla2002smote}). These have no guarantee of producing realistic samples and are weakest
in the extreme regions of the feature space. Deep generative models such as GANs~\cite{goodfellow2014gan}
and their tabular variants~\cite{xu2019ctgan} can produce richer synthetic data, but are trained to match
the bulk of the distribution and tend to under-represent its tails.

Extreme value theory (EVT)~\cite{coles2001} gives an asymptotically justified model for the tail of a
distribution. The objective of this project is a general framework that integrates EVT with deep
generative models so that synthetic minority data are statistically consistent with real extremes. The
core ideas are:
\begin{enumerate}
  \item treat the minority class as a tail phenomenon;
  \item use EVT to characterise and constrain that tail;
  \item use generative models to augment and rebalance the training data;
  \item verify that synthetic data reproduce the extreme behaviour.
\end{enumerate}
This report covers the first, basic version of the framework and its application to next-day
temperature extremes in Himachal Pradesh, India.

% -----------------------------------------------------------------------------
\section{Data}
\subsection{Sources}
Table~\ref{tab:sources} lists the data used. Raw files are never modified; every processing step
reads earlier outputs and writes to its own folder.

\begin{table}[h]
\centering\small
\caption{Data sources.}\label{tab:sources}
\begin{tabular}{@{}l p{5.2cm} l l@{}}
\toprule
Source & Content & Period, grid & Role \\
\midrule
ERA5-Land~\cite{munoz2021era5land} & Tmax, Tmin, Tmean, dewpoint, pressure, wind, soil moisture, snow,
  rain, solar radiation & 1981--2025, \ang{0.1} & labels, features \\
SRTM & elevation & \SI{30}{m} & terrain features \\
GADM boundary & state outline & -- & mask (529 cells) \\
IMD gridded temperature & Tmax, Tmin from stations & 1981--2025, \ang{1} & regional check only \\
GSOD stations & Shimla, Sundernagar, Bhuntar & 2016--2025 & planned point check \\
\bottomrule
\end{tabular}
\end{table}

ERA5-Land is the only source that is complete for 45 years at a resolution fine enough for
mountainous terrain. It has absolute biases on steep slopes, but our labels compare each grid cell with
its own history, so a constant bias shifts the threshold and not the set of extreme days.

\subsection{Pre-processing and verification}
An earlier pipeline placed all inputs on ERA5-Land's \ang{0.1} grid over
\ang{30.3}--\ang{33.3}N, \ang{75.5}--\ang{79.1}E ($31\times37$ cells), shifted the accumulated
variables (rain, radiation) back by the one day by which ERA5-Land stamps them late, averaged SRTM into
cell elevation and slope, derived relative humidity, wind speed and diurnal range, and masked cells
outside the state, leaving 529 cells and 16{,}436 days.

We verified this file against the raw ERA5-Land files for five cells (522--\SI{5279}{m}) and six
years: temperatures, dewpoint, soil moisture, snow and rain match exactly after unit conversion, with no
day where $T_{\min}>T_{\max}$. Mean temperature and humidity were missing for 1981--1985 and were
refilled from the raw files (humidity recomputed with the Magnus formula, agreeing with stored values
within 0.03\,\%RH). A bias correction against MODIS produced by the earlier pipeline was not used: its
validation against IMD showed a $-15$\,\textdegree C bias that is most plausibly the IMD missing-value
code (99.9) being treated as data.

% -----------------------------------------------------------------------------
\section{Problem formulation}
\subsection{Extreme labels}
For cell $c$ and day $t$ let $T^{\max}_{c,t}$ be the daily maximum temperature. A threshold
$u_{c,g}$ is the 99th percentile of $T^{\max}$ over the days of group $g$ in the base period 1981--2010.
Day $t$ is a heat extreme if $T^{\max}_{c,t}>u_{c,g(t)}$; cold extremes use the 1st percentile of
$T^{\min}$. Three groupings are used:
\begin{itemize}
  \item \textbf{season} (main): IMD seasons, winter (Dec--Feb), pre-monsoon (Mar--May),
        monsoon (Jun--Sep), post-monsoon (Oct--Nov); each season is compared only with the same season;
  \item \textbf{whole year}: all days;
  \item \textbf{day of year}: the same calendar day $\pm 7$ days.
\end{itemize}
Using only 1981--2010 for thresholds keeps the validation and test years unseen.
Table~\ref{tab:thr} shows that thresholds fall by about 6--7\,\textdegree C per km of altitude,
consistent with the atmospheric lapse rate.

\begin{table}[h]
\centering\small
\caption{Mean heat threshold (99th percentile of Tmax, \textdegree C) by elevation band.}\label{tab:thr}
\begin{tabular}{@{}lrrrrrr@{}}
\toprule
Elevation & Winter & Pre-monsoon & Monsoon & Post-monsoon & Whole year & Cells\\
\midrule
$<$1 km & 24.5 & 39.4 & 39.3 & 30.3 & 38.8 & 95\\
1--2 km & 20.0 & 32.0 & 31.9 & 24.4 & 31.3 & 91\\
2--3 km & 13.8 & 23.7 & 24.2 & 18.9 & 23.5 & 77\\
3--4 km & 4.4 & 13.2 & 17.5 & 12.3 & 16.7 & 65\\
4--5 km & $-3.4$ & 5.3 & 13.6 & 6.8 & 12.7 & 134\\
$>$5 km & $-8.1$ & 1.0 & 11.4 & 2.9 & 10.3 & 67\\
\bottomrule
\end{tabular}
\end{table}

\subsection{Prediction task and features}
Each sample is one cell on day $t$; the target is whether day $t+1$ is a seasonal heat extreme
(\texttt{y\_heat\_season}). Features are the weather on days $t$, $t-1$, $t-2$ (Tmax, Tmin, Tmean,
dewpoint, humidity, pressure, solar radiation, soil moisture, rain), temperature anomalies, terrain
(latitude, longitude, elevation, slope) and, from Step 3b on, distances to tomorrow's threshold
(Section~\ref{sec:features}). The data are split by year, because neighbouring days and cells are
strongly correlated (one heatwave makes hundreds of cells extreme at once):

\begin{table}[h]
\centering\small
\caption{Splits and share of extreme days.}\label{tab:splits}
\begin{tabular}{@{}llrrrr@{}}
\toprule
Split & Label years & Rows & Heat (season) & Heat (whole year) & Cold (season)\\
\midrule
Train & 1981--2010 & 5{,}794{,}666 & 1.02\% & 1.00\% & 1.02\%\\
Validation & 2011--2015 & 965{,}954 & 1.00\% & 0.77\% & 0.88\%\\
Test & 2016--2025 & 1{,}932{,}437 & 1.94\% & 1.34\% & 0.58\%\\
\bottomrule
\end{tabular}
\end{table}

\subsection{What the labels show}
\begin{itemize}
  \item \textbf{Warming.} Heat extremes nearly double in the test decade while cold extremes fall.
        2024 has the most widespread heat (4.4\% of cell-days, whole-year label).
  \item \textbf{Strongest signal in October--November.} Post-monsoon heat extremes reach 11.0\% of
        days in 2020 and 6--7\% in 2017, 2021 and 2024, against a 1\% baseline.
  \item \textbf{Clustering.} On 6 October 1983, 522 of 529 cells were extreme together; the
        59{,}000 training positives come from a few hundred events.
  \item \textbf{Season-edge effect.} Seasonal extremes gather at the warm or cold edge of a season
        (pre-monsoon heat in May, post-monsoon heat in early October, monsoon cold in late
        September); the day-of-year label serves as a comparison free of this effect.
\end{itemize}

% -----------------------------------------------------------------------------
\section{Extreme value model of the minority class}
\subsection{Peaks over threshold}
By the Pickands--Balkema--de Haan theorem~\cite{pickands1975,balkema1974}, for a high threshold $u$
the excess $Y=X-u$ given $X>u$ is approximately Generalized Pareto distributed:
\begin{equation}
  F_{\xi,\sigma}(y)=1-\left(1+\frac{\xi y}{\sigma}\right)^{-1/\xi},\qquad y>0,\; 1+\xi y/\sigma>0 .
\end{equation}
The shape $\xi$ determines the tail: $\xi<0$ gives a bounded tail with upper end point $\sigma/|\xi|$
(expected for temperature), $\xi=0$ an exponential tail and $\xi>0$ a heavy tail. Parameters are fitted
by maximum likelihood on 1981--2010 exceedances. Confidence intervals come from a block bootstrap that
resamples whole days (100 resamples), since exceedances are spatially clustered. Threshold stability is
checked by refitting above the 25th and 50th percentile of the excesses. The $m$-year return level of
the excess, for a group with $n_y$ days per year and exceedance rate $\zeta$, is
\begin{equation}
  y_m=\frac{\sigma}{\xi}\left[(m\,n_y\,\zeta)^{\xi}-1\right].
\end{equation}

\subsection{Standardised excesses}
Pooling raw excesses from all 529 cells mixes low cells with small day-to-day variability and high
cells with large variability; the mixture looks heavy-tailed even when each cell's tail is bounded. We
therefore standardise each excess by the cell's own variability $s_{c,g}$ (the standard deviation of its
daily temperature anomaly in that season, 1981--2010):
\begin{equation}
  z_{c,t}=\frac{T^{\max}_{c,t}-u_{c,g(t)}}{s_{c,g(t)}} .
\end{equation}

\begin{table}[h]
\centering\small
\caption{GPD shape $\xi$ fitted to raw and standardised excesses (1981--2010), with 95\% block-bootstrap
interval, and the share of 2016--2025 extremes beyond the fitted 90th percentile (0.10 means an
unchanged tail).}\label{tab:gpd}
\begin{tabular}{@{}lrrrr@{}}
\toprule
Group & $\xi$ raw & $\xi$ standardised & 95\% interval & Recent beyond q90\\
\midrule
Heat, whole year & $-0.153$ & $-0.152$ & $[-0.196,-0.113]$ & 0.077\\
Heat, monsoon & $-0.128$ & $-0.115$ & $[-0.185,-0.040]$ & 0.080\\
Heat, pre-monsoon & $-0.165$ & $-0.141$ & $[-0.241,-0.072]$ & 0.103\\
Heat, winter & $-0.074$ & $-0.096$ & $[-0.176,-0.032]$ & \textbf{0.209}\\
Heat, post-monsoon & $0.262$ & $0.070$ & $[-0.177,0.231]$ & 0.063\\
Heat, day of year & $0.056$ & $0.025$ & $[-0.040,0.092]$ & 0.176\\
\midrule
Cold, monsoon & $0.159$ & $-0.071$ & $[-0.142,0.023]$ & 0.005\\
Cold, pre-monsoon & $-0.111$ & $-0.061$ & $[-0.095,-0.024]$ & 0.092\\
Cold, post-monsoon & $0.085$ & $0.051$ & $[-0.021,0.112]$ & 0.149\\
Cold, winter & $0.125$ & $0.190$ & $[0.108,0.265]$ & 0.091\\
Cold, whole year & $0.010$ & $0.079$ & $[0.032,0.120]$ & 0.095\\
\bottomrule
\end{tabular}
\end{table}

After standardisation every heat tail is bounded or close to it, in line with physical expectations
(Table~\ref{tab:gpd}, Figure~\ref{fig:qq}); the spurious heavy tail of post-monsoon heat
($\xi=0.26$) largely disappears. Cold tails are only partly fixed and appear to mix different
regimes. Recent winter warm extremes exceed the old tail's 90th percentile twice as often as expected,
indicating more intense, not only more frequent, winter warm spells.

\begin{figure}[h]
\centering
\optfig{figures/evt_qq_std.png}{\linewidth}
\caption{GPD Q-Q plots of raw (grey) and standardised (blue) excesses, 1981--2010.}\label{fig:qq}
\end{figure}

% -----------------------------------------------------------------------------
\section{Baselines}\label{sec:features}
All classifiers are trained on 1981--2010; the decision threshold is chosen on 2011--2015 to maximise
F1 and applied to 2016--2025. PR-AUC is the main score (a random classifier scores the positive share,
0.019); we also report precision, recall and the critical success index
$\mathrm{CSI}=\mathrm{TP}/(\mathrm{TP}+\mathrm{FP}+\mathrm{FN})$ common in forecast verification.
\begin{itemize}
  \item \textbf{Persistence} (no model): score $=T^{\max}_{c,t}-u_{c,g(t+1)}$.
  \item \textbf{LightGBM}~\cite{ke2017lightgbm} without rebalancing, with class weights, and with SMOTE
        (minority grown to 10\% of the majority, $+514{,}294$ samples).
\end{itemize}
In a first version the model relied on calendar features (39\% of split gain), which do not transfer
to a warmer test decade. Adding the distance of today, yesterday and the day before to tomorrow's
threshold, in \textdegree C and in units of $s_{c,g}$, fixed this (the standardised distance now carries
78\% of the gain).

\begin{table}[h]
\centering\small
\caption{Test results (2016--2025) for next-day seasonal heat extremes.}\label{tab:base}
\begin{tabular}{@{}lrrrrr@{}}
\toprule
Model & PR-AUC (v1) & PR-AUC (v2) & Recall & Precision & CSI\\
\midrule
Persistence & 0.586 & 0.586 & 0.602 & 0.553 & 0.405\\
LightGBM, plain & 0.434 & 0.531 & 0.411 & 0.592 & 0.320\\
LightGBM, no calendar & -- & 0.549 & 0.468 & 0.582 & 0.350\\
LightGBM, weighted & 0.446 & 0.525 & 0.491 & 0.601 & 0.370\\
LightGBM, SMOTE & 0.449 & \textbf{0.562} & 0.515 & 0.574 & 0.373\\
\bottomrule
\end{tabular}
\end{table}

SMOTE is the strongest classic rebalancing method and is the bar for the generative methods. Persistence
remains slightly ahead overall because heat extremes persist from day to day and because the models,
trained on 1981--2010, face a test decade with twice as many extremes. In the monsoon season SMOTE
already beats persistence (PR-AUC 0.481 vs 0.442).

% -----------------------------------------------------------------------------
\section{EVT-guided CTGAN}
\subsection{Generator}
A CTGAN~\cite{xu2019ctgan}, a conditional GAN for mixed tabular data with mode-specific normalisation,
is trained on the real extreme days of a random half of 1981--2010 (about 30{,}000 samples): the 41
classifier features plus the season and the standardised severity $z$. The severity steers generation
only and is never a classifier feature, as it encodes the next day's outcome.

\subsection{Three synthetic sets from one generator}
All three come from the same trained CTGAN, so differences are due to the EVT guidance, and each adds
minority samples up to 10\% of the majority class, as for SMOTE.
\begin{enumerate}
  \item \textbf{Vanilla}: samples used as generated.
  \item \textbf{Physical}: samples violating physical consistency are removed
        ($T^{\min}\le T^{\text{mean}}\le T^{\max}$, $0\le\mathrm{RH}\le100$, non-negative rain and
        radiation, dewpoint not above temperature, consistent signs of distance features, terrain within
        the observed range).
  \item \textbf{EVT-guided} (proposed): physical filter, then for each season $k$ the samples are
        resampled so that $z$ follows the fitted $\mathrm{GPD}(\xi_k,\sigma_k)$. The support is split
        into $B=20$ bins of equal GPD probability; a sample in bin $b$ receives weight
        \begin{equation}
          w_b=\frac{1/B}{\hat p_b},\qquad \hat p_b=\frac{\#\{\text{generated } z\in b\}}{n},
        \end{equation}
        and samples with $z\ge\sigma_k/|\xi_k|$ (beyond the bounded tail) are rejected.
\end{enumerate}

\subsection{Evaluation}
The classifier is identical in every experiment (LightGBM, learning rate 0.05, up to 1{,}000 trees with
early stopping, no calendar features, and a monotone constraint that a larger distance to the threshold
can only increase the predicted risk). Besides the scores of Table~\ref{tab:base} we report
\begin{itemize}
  \item \textbf{onset} PR-AUC and recall on days that are not yet extreme, where persistence cannot help;
  \item \textbf{top-10\% recall}: the share of the most severe test extremes that are detected;
  \item \textbf{tail consistency} of each synthetic set: GPD refit on its $z$, upper quantiles and the
        Kolmogorov--Smirnov distance to the fitted tail.
\end{itemize}

\begin{table}[h]
\centering\small
\caption{Test results (2016--2025), same classifier, different training data. \pending}\label{tab:gan}
\begin{tabular}{@{}lrrrrr@{}}
\toprule
Training data & PR-AUC & Recall & CSI & Onset PR-AUC & Top-10\% recall\\
\midrule
Persistence (no model) & \pending & \pending & \pending & \pending & \pending\\
Plain & \pending & \pending & \pending & \pending & \pending\\
Class weights & \pending & \pending & \pending & \pending & \pending\\
SMOTE & \pending & \pending & \pending & \pending & \pending\\
CTGAN, vanilla & \pending & \pending & \pending & \pending & \pending\\
CTGAN, physical filter & \pending & \pending & \pending & \pending & \pending\\
\textbf{CTGAN, EVT-guided} & \pending & \pending & \pending & \pending & \pending\\
\bottomrule
\end{tabular}
\end{table}

\begin{figure}[h]
\centering
\optfig{figures/tail_survival.png}{\linewidth}
\caption{Survival curves of the standardised severity $z$ per season: real extremes, vanilla CTGAN,
EVT-guided CTGAN and the fitted GPD.}\label{fig:tail}
\end{figure}

\begin{figure}[h]
\centering
\optfig{figures/pr_curves.png}{0.6\linewidth}
\caption{Precision--recall curves on the test decade 2016--2025.}\label{fig:pr}
\end{figure}

% -----------------------------------------------------------------------------
\section{Discussion and limitations}
\begin{itemize}
  \item Labels come from a reanalysis; a check against station observations (GSOD: Shimla,
        Sundernagar, Bhuntar) is planned.
  \item The test decade has about twice the training rate of heat extremes. This non-stationarity
        hurts any model trained on 1981--2010 and is itself a reason to model the tail explicitly.
  \item Seasonal labels show a season-edge effect for post-monsoon heat and monsoon cold.
  \item Cold tails are not yet captured by a single pooled GPD.
  \item EVT currently guides the generator only at sampling time; it does not shape training.
  \item All experiments run on free Colab resources.
\end{itemize}

\section{Next steps}
\begin{enumerate}
  \item Complete and analyse the Step 4 comparison (Table~\ref{tab:gan}).
  \item Move EVT into training: a GPD-likelihood penalty on generated excesses in the generator loss.
  \item Apply the same guidance to a tabular diffusion model (TabDDPM~\cite{kotelnikov2023tabddpm}).
  \item Extend to cold extremes, add station checks, and test on public imbalanced benchmarks to
        support the general-purpose claim.
\end{enumerate}

\section*{Reproducibility}
All code is in the GitHub repository \texttt{Divyansh-jain2/Major\_Technical\_Project}
(folder \texttt{steps/}). Each step runs as one Colab cell, reads only earlier outputs, writes only to
its own folder and saves a text summary.

% -----------------------------------------------------------------------------
\begin{thebibliography}{99}
\bibitem{coles2001} S.~Coles. \textit{An Introduction to Statistical Modeling of Extreme Values}.
  Springer, 2001.
\bibitem{pickands1975} J.~Pickands. Statistical inference using extreme order statistics.
  \textit{Annals of Statistics}, 3(1):119--131, 1975.
\bibitem{balkema1974} A.~A.~Balkema and L.~de~Haan. Residual life time at great age.
  \textit{Annals of Probability}, 2(5):792--804, 1974.
\bibitem{chawla2002smote} N.~V.~Chawla, K.~W.~Bowyer, L.~O.~Hall and W.~P.~Kegelmeyer. SMOTE: Synthetic
  minority over-sampling technique. \textit{Journal of Artificial Intelligence Research}, 16:321--357, 2002.
\bibitem{goodfellow2014gan} I.~Goodfellow et al. Generative adversarial nets. In \textit{Advances in
  Neural Information Processing Systems}, 2014.
\bibitem{xu2019ctgan} L.~Xu, M.~Skoularidou, A.~Cuesta-Infante and K.~Veeramachaneni. Modeling tabular
  data using conditional GAN. In \textit{Advances in Neural Information Processing Systems}, 2019.
\bibitem{kotelnikov2023tabddpm} A.~Kotelnikov, D.~Baranchuk, I.~Rubachev and A.~Babenko. TabDDPM:
  Modelling tabular data with diffusion models. In \textit{International Conference on Machine
  Learning}, 2023.
\bibitem{ke2017lightgbm} G.~Ke et al. LightGBM: A highly efficient gradient boosting decision tree. In
  \textit{Advances in Neural Information Processing Systems}, 2017.
\bibitem{munoz2021era5land} J.~Mu\~noz-Sabater et al. ERA5-Land: a state-of-the-art global reanalysis
  dataset for land applications. \textit{Earth System Science Data}, 13:4349--4383, 2021.
\end{thebibliography}

\end{document}
